\chapter*{ВВЕДЕНИЕ}
\addcontentsline{toc}{chapter}{ВВЕДЕНИЕ}
\sloppy


\section*{Актуальность темы}
\addcontentsline{toc}{section}{Актуальность темы}

Современные научно-практические IT-конференции — многодневные мероприятия с десятками параллельных сессий, в которых участвуют тысячи специалистов. Программа конференции включает доклады разной тематики, уровня и популярности; залы существенно различаются по вместимости. Распределение участников между параллельными сессиями определяет содержательную ценность мероприятия для участников, репутацию и коммерческую привлекательность мероприятия для организаторов и охват аудитории для докладчиков.

В отсутствие координирующего вмешательства распределение формируется как агрегат независимых индивидуальных выборов. Это приводит к одновременной перегрузке популярных залов, которые не могут вместить желающих, и хронической недозагрузке менее популярных, что снижает суммарное качество мероприятия. Прямое применение классических рекомендательных систем не решает задачи: оптимизируя индивидуальную релевантность независимо для каждого пользователя, такие системы воспроизводят одинаково популярные позиции в верхних строках выдачи и тем самым усиливают, а не ослабляют нагрузку на наиболее востребованные ресурсы.

Описанная задача относится к более широкому классу задач распределения автономных пользователей по ресурсам ограниченной вместимости через ненавязчивую рекомендательную подсказку. Этот класс задач возникает в маршрутизации транспортных потоков, управлении сервисами бронирования, образовательных платформах, управлении посетителями выставочных пространств. Несмотря на разнообразие приложений, существующие подходы не закрывают совокупности требований, характерной для рассматриваемой постановки: совместная оптимизация рекомендаций между пользователями, жёсткие ограничения вместимости в одной коллективной выдаче, ограниченность исторических данных, мягкий характер управляющего рычага.

Развитие рекомендательных систем, методов обучения с подкреплением и симуляционных подходов с использованием больших языковых моделей в 2023–2026 годах расширило инструментарий, применимый к задачам этого класса. Подходы на основе генеративных агентов и LLM-симуляторов допускают работу с малыми каталогами при отсутствии исторических данных взаимодействий; методы обучения политики с ограничениями дают формальный аппарат для учёта системных ограничений вместимости. Систематический синтез этих подходов в задаче формирования программы конференции составляет содержание настоящей работы.

\section*{Степень разработанности темы}
\addcontentsline{toc}{section}{Степень разработанности темы}

Задача в её обобщённой постановке исследовалась с разных сторон. В рамках теории игр аппарат игр с перегрузкой Розенталя [Rosenthal, 1973] и равновесия Уордропа [Wardrop, 1952] обеспечивают концептуальное описание ситуации децентрализованного выбора при взаимной зависимости пользователей. В рекомендательных системах подходы ранжирования с ограничениями Сингха и Йоахимса [Singh, Joachims, 2018] и калиброванных рекомендаций Стека [Steck, 2018] вводят аппарат учёта системных свойств выдачи. В обучении с подкреплением марковский процесс принятия решений с ограничениями [Altman, 1999] и его применения в производственных рекомендательных системах [Chen et al., 2019; Zheng et al., 2018] предоставляют операциональный механизм построения управляющей политики. Современные направления использования больших языковых моделей в роли симуляторов пользователей [Park et al., 2023; Zhang et al., 2024] и в роли ранжирующих компонентов [Hou et al., 2024] открывают возможность работы с малыми каталогами и качественной валидации моделей пользовательского поведения.

Вместе с тем, систематическое объединение указанных подходов применительно к задаче рассматриваемой структуры — формирование согласованных рекомендательных выдач для множества пользователей одновременно при жёстких ограничениях вместимости в одной коллективной выдаче — в открытой литературе отсутствует. Это определяет содержательное пространство для разработки методологического и алгоритмического подхода в работе.

\section*{Объект и предмет исследования}
\addcontentsline{toc}{section}{Объект и предмет исследования}

\textbf{Объект исследования} — процесс формирования персональных программ участников научно-практических IT-конференций при ограничениях вместимости параллельных сессий.

\textbf{Предмет исследования} — методы построения управляющей рекомендательной политики, согласующей индивидуальные предпочтения участников с системными ограничениями вместимости и обеспечивающей совместную оптимизацию рекомендаций между пользователями.

\section*{Цель работы}
\addcontentsline{toc}{section}{Цель работы}

Цель работы — разработать интеллектуальную систему поддержки формирования программы конференции, учитывающую ограничения вместимости параллельных сессий и согласованность рекомендаций между участниками.

\section*{Задачи исследования}
\addcontentsline{toc}{section}{Задачи исследования}

Для достижения поставленной цели в работе решаются следующие задачи:

\begin{enumerate}
  \item Провести анализ предметной области, систематизировать существующие концептуальные рамки и методы решения задач рассматриваемого класса, выделить пробелы в существующих подходах.
  \item Формализовать постановку задачи как марковский процесс принятия решений с ограничениями в структуре игры с перегрузкой; сформулировать функциональные и нефункциональные требования к разрабатываемому решению.
  \item Разработать гибридную симуляционную среду для обучения и оценки рекомендательных политик в условиях отсутствия исторических данных взаимодействий, объединяющую однократную генерацию синтетической пользовательской популяции на основе большой языковой модели и параметрический пошаговый симулятор.
  \item Реализовать набор сравниваемых рекомендательных политик различной природы — случайную, контентную на основе косинусной близости, эвристическую с учётом вместимости, основанную на алгоритме Maximal Marginal Relevance, обучаемую методом проксимальной оптимизации в постановке марковского процесса с ограничениями и опирающуюся на большую языковую модель — и провести их сравнение по совокупности метрик качества.
  \item Провести независимую валидацию политики, выбранной по результатам основного сравнения, в среде на основе автономных агентов с большими языковыми моделями; оценить разрыв между симуляцией и реальным поведением пользователей.
\end{enumerate}

\section*{Методы исследования}
\addcontentsline{toc}{section}{Методы исследования}

Работа опирается на анализ литературы по рекомендательным системам, теории игр с перегрузкой, обучению с подкреплением и стохастической оптимизации; математическое моделирование задачи как марковского процесса принятия решений с ограничениями; разработку симуляционной среды с использованием параметрических моделей пользовательского выбора и генеративных агентов на основе больших языковых моделей; обучение политики методами обучения с подкреплением; статистическое сравнение результатов нескольких рекомендательных политик по совокупности метрик качества.

\section*{Научная новизна}
\addcontentsline{toc}{section}{Научная новизна}

Научная новизна работы заключается в систематическом синтезе элементов нескольких теоретических рамок применительно к задаче рассматриваемого класса. Конкретно: сочетание формализации задачи как марковского процесса принятия решений с ограничениями со структурным обоснованием через теорию игр с перегрузкой, использование гибридной симуляционной среды с однократной генерацией пользовательской популяции на основе большой языковой модели для решения задач с малым каталогом ресурсов, двухслойная схема валидации с независимой проверкой в среде на основе автономных агентов. Указанное сочетание адресует совокупность требований, не закрываемую ни одним из существующих подходов в отдельности.

\section*{Практическая значимость}
\addcontentsline{toc}{section}{Практическая значимость}

Практическая значимость работы определяется применимостью разработанной системы к реальной задаче формирования программы конференций — на примере конференций Mobius и Heisenbug компании JUG. Эмпирическая часть выполняется на материале открытых программ конференций Mobius и Heisenbug 2023–2025 годов (около 30–40 докладов на конференцию с описаниями тематики, форматов и распределением по тайм-слотам и залам); параметры симуляционной среды калибруются по фактической вместимости залов и типичной численности участников. Разработанные симуляционная среда и набор сравниваемых политик могут быть переиспользованы для решения задач смежных доменов, входящих в рассматриваемый класс: маршрутизация транспортных потоков, бронирование сервисов с ограниченной вместимостью, управление учебной нагрузкой в образовательных платформах, управление потоками посетителей в зонах с пропускными ограничениями.

\section*{Положения, выносимые на защиту}
\addcontentsline{toc}{section}{Положения, выносимые на защиту}

\begin{enumerate}
  \item Постановка задачи формирования согласованных рекомендательных выдач при ограничениях вместимости как марковского процесса принятия решений с ограничениями в структуре игры с перегрузкой обеспечивает формальный аппарат для совместной оптимизации рекомендаций между пользователями при жёстких ограничениях вместимости в одной коллективной выдаче.
  \item Гибридная симуляционная среда, объединяющая однократную генерацию пользовательской популяции на основе большой языковой модели с параметрическим пошаговым симулятором, обеспечивает обучаемость рекомендательной политики при отсутствии реальных данных взаимодействий и контролируемой стоимости вычислений.
  \item Сравнительное исследование политик различной природы — случайной, контентной, эвристической с учётом вместимости, основанной на алгоритме MMR, обучаемой методом проксимальной оптимизации в постановке марковского процесса с ограничениями и опирающейся на большую языковую модель — позволяет систематически охарактеризовать компромиссы между соблюдением системных ограничений вместимости и пользовательской полезностью на материале рассматриваемой задачи.
  \item Двухслойная схема валидации с проверкой выбранной политики во второй симуляционной среде на основе автономных агентов с большими языковыми моделями позволяет оценить устойчивость результатов основного сравнения к смене модели пользовательского поведения и сформировать содержательную интерпретацию расхождений между двумя моделями симуляции.
\end{enumerate}

\section*{Апробация результатов}
\addcontentsline{toc}{section}{Апробация результатов}

Результаты работы получены и проверены на разработанной двухслойной симуляционной среде, объединяющей параметрический пошаговый симулятор и независимый слой валидации на основе автономных агентов с большими языковыми моделями. Подобная схема обеспечивает воспроизводимое сравнение методов и независимый контроль внешней валидности результатов. Постановка задачи и предлагаемый методологический подход обсуждались с научным руководителем работы и представителями команды разработки рекомендательного сервиса конференций Mobius и Heisenbug компании JUG, на материале которых выполняется эмпирическая часть работы.

\section*{Структура работы}
\addcontentsline{toc}{section}{Структура работы}

Работа состоит из введения, четырёх глав, заключения, списка использованных источников и приложений. В главе 1 проводится анализ предметной области, систематизируются концептуальные рамки и методы, идентифицируется пробел в существующих подходах. В главе 2 формализуется постановка задачи, формулируются требования к решению, описывается предлагаемый метод и определяется дизайн экспериментальной части работы. В главе 3 представляются технические решения по реализации симуляционной среды и сравниваемых политик, описывается ход проведённых экспериментов. В главе 4 приводятся результаты сравнения политик, ablation studies, проверка устойчивости и итоговая валидация в среде на основе автономных агентов; обсуждается разрыв между симуляцией и реальным поведением. В заключении формулируются основные результаты работы и направления дальнейших исследований.
